\documentclass[11pt,a4paper]{article}
\usepackage[utf8]{inputenc}
\usepackage[vietnamese]{babel}
\usepackage{amsmath, amsfonts, amssymb}
\usepackage{graphicx}
\usepackage{geometry}
\usepackage{booktabs}
\usepackage{enumitem}
\usepackage{setspace}
\usepackage{fancyhdr}
\usepackage{float}
\usepackage{listings}
\usepackage{xcolor}
\usepackage{hyperref}
\usepackage[most]{tcolorbox}

% --- Cấu hình lề và khoảng cách ---
\geometry{margin=1in, top=1.2in, bottom=1in}
\onehalfspacing

% --- Cấu hình Code Snippets ---
\lstset{
    commentstyle=\color{gray},
    keywordstyle=\bfseries\color{blue},
    numberstyle=\tiny\color{gray},
    stringstyle=\color{red},
    basicstyle=\ttfamily\footnotesize,
    breaklines=true,
    frame=single,
    tabsize=2
}

% --- Cấu hình Header/Footer ---
\pagestyle{fancy}
\fancyhf{}
\lhead{Đồ án Mã hóa ứng dụng và an ninh thông tin}
\rhead{Giai đoạn 6: Bảo trì Phần mềm}
\cfoot{\thepage}

\begin{document}

% --- TRANG TIÊU ĐỀ ---
\begin{center}
    {\Large \textbf{BÁO CÁO GIAI ĐOẠN CUỐI}} \\
    \vspace{0.5em}
    {\huge \textbf{BẢO TRÌ VÀ PHÁT TRIỂN HỆ THỐNG}} \\
    \vspace{0.5em}
    \textbf{Tên sản phẩm: SEC-Wallet (Security Core)}
\end{center}

\noindent
\begin{tabular}{@{}l l p{3cm} l}
\textbf{Sinh viên thực hiện:} & Thái Hữu Thọ & & \textbf{Lớp:} [Tên lớp] \\
\textbf{MSSV:} & 22127400 & & \textbf{Ngày:} \today \\
\end{tabular}

\vspace{0.5cm}
\rule{\linewidth}{1pt}
\vspace{0.8cm}

\section{Tổng quan về Giai đoạn Bảo trì (Software Maintenance)}
Trong vòng đời phát triển phần mềm (SDLC), giai đoạn bảo trì không chỉ đơn thuần là sửa lỗi mà còn là quá trình tiến hóa để hệ thống thích nghi với môi trường mới và nâng cao trải nghiệm người dùng. Đối với dự án **SEC-Wallet**, bảo trì đóng vai trò sống còn vì các lỗ hổng mật mã và lỗi logic giao dịch có thể dẫn đến mất mát tài sản thực tế.

Báo cáo này liệt kê chi tiết các hoạt động bảo trì đã được thực hiện dựa trên 4 loại hình chuẩn: Corrective (Sửa chữa), Adaptive (Thích nghi), Perfective (Hoàn thiện) và Preventive (Phòng ngừa).

\section{Bảo trì Sửa chữa (Corrective Maintenance)}
Đây là các hoạt động tập trung vào việc khắc phục các lỗi logic và kỹ thuật phát sinh trong quá trình vận hành ban đầu.

\subsection{1. Khắc phục lỗi sai lệch Schema (Database Mapping Fix)}
Trong quá trình phát triển tính năng "Yêu cầu thanh toán" (Payment Request), hệ thống đã gặp lỗi 500 do sự không tương thích giữa tham số đầu vào và mô hình cơ sở dữ liệu.
\begin{itemize}
    \item \textbf{Vấn đề:} Model \texttt{Transaction} yêu cầu trường \texttt{aad} (Additional Authenticated Data), trong khi logic khởi tạo lại gửi vào trường \texttt{message}.
    \item \textbf{Giải pháp:} Ánh xạ lại toàn bộ luồng dữ liệu, đóng gói \texttt{tx\_id} và các metadata khác vào khối \texttt{aad} theo tiêu chuẩn AES-GCM.
\end{itemize}

\subsection{2. Sửa lỗi logic xác minh mã PIN (Verification Logic Fix)}
Một lỗi nghiêm trọng đã được phát hiện khi hệ thống sử dụng băm mật khẩu đăng nhập (\texttt{password\_hash}) để kiểm tra mã PIN thanh toán thay vì sử dụng \texttt{payment\_pin\_hash}.
\begin{itemize}
    \item \textbf{Giải pháp:} Cập nhật endpoint \texttt{/fulfill} để truy vấn chính xác trường mã PIN đã được băm bằng Argon2, đảm bảo tính tách biệt giữa mật khẩu truy cập và mã PIN bảo vệ dòng tiền.
\end{itemize}

\subsection{3. Cưỡng chế các ràng buộc nghiệp vụ (Business Rule Enforcement)}
Bổ sung các lớp kiểm tra logic để ngăn chặn các giao dịch phi lý:
\begin{itemize}
    \item Chặn hoàn toàn khả năng người dùng tự gửi yêu cầu thanh toán cho chính mình.
    \item Chặn giao dịch đến tài khoản \texttt{admin} (tài khoản hệ thống chỉ dùng để quản trị).
\end{itemize}

\section{Bảo trì Thích nghi (Adaptive Maintenance)}
Thích nghi hệ thống với những thay đổi về hạ tầng và công nghệ.

\subsection{1. Chuyển đổi kiến trúc từ Desktop sang Full-stack Web}
Ban đầu hệ thống được thiết kế dạng ứng dụng Desktop Local. Để mở rộng khả năng tiếp cận và demo báo cáo, hệ thống đã được chuyển đổi sang kiến trúc Web:
\begin{itemize}
    \item \textbf{Backend:} Sử dụng FastAPI với cơ chế xử lý bất đồng bộ (Asynchronous).
    \item \textbf{Persistence:} Chuyển từ lưu trữ file JSON đơn giản sang SQLite (\texttt{sec\_wallet.db}) để đảm bảo tính Acid (Atomicity, Consistency, Isolation, Durability) cho giao dịch.
\end{itemize}

\subsection{2. Cơ chế Cache-Busting cho Frontend}
Để đảm bảo các bản sửa lỗi JavaScript được cập nhật ngay lập tức trên trình duyệt của người chấm bài mà không cần xóa cache thủ công, hệ thống đã áp dụng kỹ thuật versioning:
\begin{lstlisting}[language=html]
<script src="script.js?v=4.0"></script>
\end{lstlisting}

\section{Bảo trì Hoàn thiện (Perfective Maintenance)}
Nâng cao hiệu năng và trải nghiệm người dùng (UX) dựa trên phản hồi thực tế.

\subsection{1. Đồng bộ hóa trải nghiệm người dùng (UX Standardization)}
Giao diện "Yêu cầu tiền" (Request Money) đã được tái thiết kế hoàn toàn để khớp với quy trình "Gửi tiền" (Send Money):
\begin{itemize}
    \item Sử dụng luồng: \textbf{Nhập liệu $\rightarrow$ Hiện Modal bảo mật $\rightarrow$ Xác nhận mã PIN}.
    \item Tự động xóa trắng form (Input resets) khi chuyển tab hoặc sau khi thực hiện giao dịch thành công để nâng cao tính riêng tư.
\end{itemize}

\subsection{2. Làm sạch mã nguồn và Tái cấu trúc (Code Refactoring)}
Để giảm nợ kỹ thuật (Technical Debt), các module bảo mật đã được tách biệt rõ ràng:
\begin{itemize}
    \item \texttt{encryptor.py}: Chuyên trách AES-GCM và EdDSA Signing.
    \item \texttt{decryptor.py}: Chuyên trách Verify và Decrypt.
    \item \texttt{keygen.py}: Chuyên trách quản lý cặp khóa ECC.
\end{itemize}

\section{Bảo trì Phòng ngừa (Preventive Maintenance)}
Các biện pháp ngăn chặn sự cố trước khi chúng xảy ra.

\subsection{1. Nâng cấp tham số an ninh (Security Hardening)}
Dựa trên khuyến nghị của NIST, hệ thống đã nâng cấp số vòng lặp PBKDF2 từ 100,000 lên **600,000 vòng** cho việc mã hóa file khóa private key, giúp kéo dài thời gian tấn công vét cạn (Brute-force) của kẻ gian.

\subsection{2. Giới hạn tần suất và Khóa tài khoản (Account Lockout)}
Triển khai thư viện \texttt{slowapi} để giới hạn số lần thử đăng nhập. Nếu sai quá 5 lần, tài khoản sẽ tự động khóa trong 15 phút, ngăn chặn các cuộc tấn công Brute-force tự động.

\section{Kết luận}
Quá trình bảo trì Giai đoạn 5 đã giúp **SEC-Wallet** trở nên ổn định hơn, an toàn hơn và chuyên nghiệp hơn. Việc xử lý triệt để các lỗi logic giao dịch và nâng cấp kiến trúc lên Web đã chứng minh khả năng tiến hóa của phần mềm để đáp ứng các tiêu chuẩn khắt khe trong ngành an ninh thông tin.

\end{document}
